\documentclass[12pt, a4paper]{article}

% --- Encoding & language ---
\usepackage[utf8]{inputenc}
\usepackage[T1]{fontenc}
\usepackage[english]{babel}

% --- Layout ---
\usepackage[margin=2.5cm]{geometry}
\usepackage{parskip}
\usepackage{titlesec}
\usepackage{enumitem}

% --- Graphics & colors ---
\usepackage{graphicx}
\usepackage{xcolor}

% --- Math ---
\usepackage{amsmath}
\usepackage{amssymb}

% --- Tables ---
\usepackage{booktabs}
\usepackage{tabularx}
\usepackage{array}
\usepackage{longtable}

% --- Code listings ---
\usepackage{listings}

\definecolor{codebg}{HTML}{F5F5F5}
\definecolor{codeframe}{HTML}{CCCCCC}
\definecolor{keywordcolor}{HTML}{0000AA}
\definecolor{commentcolor}{HTML}{666666}
\definecolor{stringcolor}{HTML}{AA0000}
\definecolor{prologatom}{HTML}{007700}

\lstdefinelanguage{Prolog}{
  keywords={is, not, true, fail, assert, retract, asserta, assertz,
            retractall, write, writeln, nl, if, random, findall,
            forall, format, atom_concat, number_codes, atom_codes},
  keywordstyle=\color{keywordcolor}\bfseries,
  comment=[l]{\%},
  commentstyle=\color{commentcolor}\itshape,
  stringstyle=\color{stringcolor},
  morestring=[b]',
  morestring=[b]",
  sensitive=false
}

\lstset{
  basicstyle=\ttfamily\small,
  backgroundcolor=\color{codebg},
  frame=single,
  rulecolor=\color{codeframe},
  breaklines=true,
  breakatwhitespace=false,
  numbers=left,
  numberstyle=\tiny\color{gray},
  numbersep=8pt,
  xleftmargin=15pt,
  keywordstyle=\color{keywordcolor},
  commentstyle=\color{commentcolor}\itshape,
  stringstyle=\color{stringcolor},
  showstringspaces=false,
  tabsize=2
}

% --- Hyperlinks ---
\usepackage[hidelinks, colorlinks=false]{hyperref}

% ----------------------------------------------------------------
% Title info
% ----------------------------------------------------------------
\title{%
  \vspace{-1.5cm}
  {\Large F1 Race}\\[0.4em]
  {\large DALI Multi-Agent System --- Technical Documentation}
}
\author{Piccirilli Michael}
\date{\today}

% ================================================================
\begin{document}

\maketitle

\tableofcontents
\newpage

\begin{abstract}
This document describes the design, architecture, and implementation of a Formula~1 race
simulation built on top of the DALI Multi-Agent System framework. Car agents are
dynamically generated from a single agents.json configuration file --- no
agents are hardcoded in the codebase. A dedicated Pit-Wall agent coordinates the
entire race flow and injects probabilistic track events (safety car deployment, heavy rain)
at runtime. A real-time web dashboard provides live monitoring and manual control of the
simulation.
\end{abstract}

\bigskip

% ================================================================
\section{Introduction}

\subsection{DALI Overview}
DALI (Dynamic Agent Logic framework for Interaction) is a logic-based multi-agent system
built on top of SICStus Prolog. Each agent is defined by a set of \emph{external events}
(reactions to incoming messages), \emph{internal events} (proactive behaviours triggered by
conditions), and standard Prolog clauses for knowledge representation and reasoning.

DALI agents run as independent SICStus processes that communicate via a shared message
server running on a configurable TCP port (default: \texttt{3010}).

\subsection{Project Goals}
The F1 Race simulation was developed with the following goals in mind:
\begin{itemize}[leftmargin=2em]
  \item Demonstrate the \emph{reactive} and \emph{proactive} capabilities of DALI agents in
        a realistic, event-driven scenario.
  \item Show how a fully \emph{dynamic} agent roster can be maintained through code
        generation, avoiding any hardcoded agent names in the DALI source.
  \item Provide a real-time monitoring interface that reflects agent activity without
        modifying the agent logic.
  \item Package the entire stack as a reproducible \textbf{Docker Compose} application.
\end{itemize}

% ================================================================
\section{System Requirements}

\begin{longtable}{@{} l l @{}}
\toprule
\textbf{Component} & \textbf{Requirement} \\
\midrule
\endhead
SICStus Prolog & 4.6.0 at \texttt{/usr/local/sicstus4.6.0} (Linux/WSL) \\
Shell           & \texttt{bash}, \texttt{tmux} \\
Python          & 3.8+ (for \texttt{generate\_agents.py} and dashboard) \\
Flask           & Auto-installed in a local venv by \texttt{ui/run.sh} \\
Docker          & Optional; required only for the Docker-based deployment \\
\bottomrule
\end{longtable}

% ================================================================
\section{Project Structure}

\begin{lstlisting}[language={}, numbers=none, frame=none, backgroundcolor=\color{white}]
f1_race/
+-- agents.json           <- Single source of truth for car agents
+-- generate_agents.py    <- Generates DALI files from agents.json
+-- startmas.sh           <- Launch script (Linux/WSL)
+-- startmas.bat          <- Launch script (Windows)
+-- docker-compose.yml    <- Docker Compose (mas + ui containers)
+-- .env.example          <- Template for SICSTUS_PATH
+-- docker/
|   +-- setup.sh          <- Auto-detects SICStus, writes .env
|   +-- mas/Dockerfile
|   \-- ui/Dockerfile
+-- ui/
|   +-- dashboard.py      <- Flask backend
|   +-- run.sh            <- Creates venv + launches dashboard
|   +-- requirements.txt
|   \-- static/
|       +-- app.js        <- Frontend (syncConfig every 5 s)
|       \-- index.html
+-- mas/
|   +-- instances/        <- Agent instance files
|   \-- types/            <- Agent type (behaviour) files
+-- conf/                 <- DALI communication configuration
+-- build/                <- Runtime binaries (auto-generated)
+-- work/                 <- Runtime working files (auto-generated)
\-- log/                  <- Agent logs (auto-generated)
\end{lstlisting}

% ================================================================
\section{Agents}

\subsection{Fixed Infrastructure Agents}

The following agents are always present regardless of the car roster:

\begin{longtable}{@{} l l l p{7cm} @{}}
\toprule
\textbf{Agent} & \textbf{Instance file} & \textbf{Type} & \textbf{Role} \\
\midrule
\endhead
\texttt{semaphore}   & \texttt{instances/semaphore.txt}   & \texttt{semaphoreType}  & Collects \texttt{ready} signals, runs the F1 lights sequence, fires \texttt{start\_race} \\[4pt]
\texttt{pitwall}     & \texttt{instances/pitwall.txt}     & \texttt{pitWallType}    & Race coordinator: manages lap order, lap times, standings, track events \\[4pt]
\texttt{safety\_car} & \texttt{instances/safety\_car.txt} & \texttt{safetyCarType}  & Reacts to deploy/recall; forwards \texttt{sc\_recalled} and participates in the green-flag chain \\
\bottomrule
\end{longtable}

\subsection{Car Agents (Dynamic)}

Car agents are \emph{never hardcoded}. The entire roster is driven by
\texttt{agents.json}. The default configuration defines four cars:

\begin{longtable}{@{} l l l l @{}}
\toprule
\textbf{ID} & \textbf{Team} & \textbf{Driver} & \textbf{Car model} \\
\midrule
\endhead
\texttt{ferrari}  & Ferrari  & Leclerc    & SF-24 \\
\texttt{mclaren}  & McLaren  & Norris     & MCL38 \\
\texttt{redbull}  & Red Bull & Verstappen & RB20  \\
\texttt{mercedes} & Mercedes & Hamilton   & W15   \\
\bottomrule
\end{longtable}

\noindent
Each car entry in \texttt{agents.json} also carries UI metadata (\texttt{color},
\texttt{border}, \texttt{label}) used by the dashboard.

% ================================================================
\section{Dynamic Agent System}

\subsection{Code Generation Pipeline}

The \texttt{generate\_agents.py} script is the single point of truth for all generated
files:

\begin{lstlisting}[language={}, numbers=none, frame=none, backgroundcolor=\color{white}]
agents.json
    +-> generate_agents.py
          +-- mas/instances/{id}.txt       (one per car)
          +-- mas/types/{id}Car.txt        (one per car)
          +-- mas/types/pitWallType.txt    (round-robin over all cars)
          +-- mas/types/semaphoreType.txt  (waits N_cars + 2 ready)
          \-- mas/types/safetyCarType.txt
\end{lstlisting}

\noindent
\texttt{generate\_agents.py} is invoked automatically at each MAS startup by
\texttt{startmas.sh}. It is \emph{smart about regeneration}:

\begin{longtable}{@{} l p{9cm} @{}}
\toprule
\textbf{Situation} & \textbf{Behaviour} \\
\midrule
\endhead
\texttt{agents.json} unchanged & Skips regeneration --- no files written \\
New car added                  & Full regeneration \\
Car removed                    & Stale \texttt{\{id\}.txt} / \texttt{\{id\}Car.txt} deleted, then full regeneration \\
\texttt{-{}-force} flag        & Always regenerates unconditionally \\
\bottomrule
\end{longtable}

\subsection{Adding or Removing a Car}

Edit \texttt{agents.json} and add a new entry to the \texttt{cars} array:

\begin{lstlisting}[language={}]
{
  "total_laps": 5,
  "cars": [
    { "id": "ferrari",   "team": "Ferrari",  "car_model": "SF-24",
      "driver": "Leclerc", "label": "Ferrari SF-24",
      "color": "#180505", "border": "#cc2200" },
    { "id": "myclubcar", "team": "My Club",  "car_model": "X1",
      "driver": "Rossi",   "label": "Club X1",
      "color": "#001020", "border": "#00aaff" }
  ]
}
\end{lstlisting}

\noindent
Then run \texttt{bash startmas.sh}. All DALI files and the dashboard update automatically.

% ================================================================
\section{Race Flow}

\subsection{Startup Sequence}

\begin{enumerate}[leftmargin=2em]
  \item Every agent (all cars, pit wall, safety car) sends a \texttt{ready} message to the
        semaphore.
  \item The semaphore waits until it has collected exactly $N_\text{cars} + 2$ ready signals.
  \item The semaphore runs the F1 lights sequence (5 lights on, 1\,s each; 2\,s pause;
        lights out) and sends \texttt{start\_race} to \texttt{Car[0]}.
\end{enumerate}

\subsection{Lap Round-Robin}

Once the race starts, laps proceed in a strict round-robin over the ordered car list:

\begin{enumerate}[leftmargin=2em]
  \item \texttt{Car[i]} receives \texttt{lap\_go\_\{id\}} from the pit wall.
  \item \texttt{Car[i]} asserts its lap status and sends \texttt{lap\_done\_\{id\}} to the
        pit wall.
  \item The pit wall adds a random lap time, optionally fires a track event, and sends
        \texttt{lap\_go\_\{id+1\}} to the next car.
  \item After the last car of a lap completes:
        \begin{itemize}
          \item Lap counter increments.
          \item Standings are printed.
          \item If laps remain: \texttt{random\_track\_event} is rolled, then \texttt{lap\_go}
                is sent to \texttt{Car[0]}.
          \item If race complete: \texttt{declare\_winner} is called and \texttt{race\_end}
                is broadcast.
        \end{itemize}
\end{enumerate}

\noindent
At any point a car's \emph{internal events} (\texttt{engine\_failureI} or
\texttt{push\_lapI}) can fire autonomously and notify the pit wall.

% ================================================================
\section{Timing System}

The winner is the car with the lowest total accumulated time. Time is measured in seconds.

\begin{longtable}{@{} l r @{}}
\toprule
\textbf{Event} & \textbf{Time change} \\
\midrule
\endhead
Each lap      & $+\,\texttt{random}(60,\,90)$ seconds \\
Pit stop      & $+\,25\,\text{s}$ \\
Safety car (20\,\% per lap) & $+\,10\,\text{s}$ to \emph{all} cars \\
Heavy rain (20\,\% per lap) & $+\,5\,\text{s}$ to \emph{all} cars \\
Push lap (internal event, $\approx$10\,\%) & $-\,3\,\text{s}$ \\
Engine failure / DNF ($\approx$0.2\,\%) & $\texttt{time} = 9999\,\text{s}$ (race ends) \\
\bottomrule
\end{longtable}

% ================================================================
\section{DALI Event Types}

\subsection{External Events}

External events (\texttt{nameE:>~Body.}) are reactive: they fire when a message named
\texttt{name} arrives from another agent.

\begin{longtable}{@{} l p{9cm} @{}}
\toprule
\textbf{Event} & \textbf{Description} \\
\midrule
\endhead
\texttt{start\_raceE}             & Lights-out signal from semaphore; asserts \texttt{race\_started} \\
\texttt{lap\_go\_\{id\}E}        & Pit wall tells car to start its next lap \\
\texttt{lap\_done\_\{id\}E}      & Car finishes a lap; pit wall adds time, rolls track event \\
\texttt{pit\_done\_\{id\}E}      & Car finishes pit stop; pit wall adds $+25\,\text{s}$ \\
\texttt{\{id\}\_engine\_failureE} & DNF notification; sets effective time to 9999 \\
\texttt{\{id\}\_push\_lapE}      & Fastest-lap bonus; subtracts $3\,\text{s}$ \\
\texttt{race\_endE}              & Broadcast at end of race; asserts local \texttt{race\_over} \\
\texttt{rain\_warningE}          & Cosmetic notification (ignored if \texttt{race\_over}) \\
\texttt{safety\_car\_deployedE}  & Car receives safety-car notification \\
\texttt{green\_flagE}            & Sent after safety car recalled; car reacts with a push message \\
\texttt{retire\_\{id\}E}         & Car parks; cosmetic DNF notification \\
\bottomrule
\end{longtable}

\subsection{Internal Events}

Internal events (\texttt{nameI:>~Body.}) are proactive: they fire autonomously whenever
the associated condition clause \texttt{name~:-~Cond.} succeeds.

\begin{lstlisting}[language=Prolog, caption={Generated internal events (per-car)}]
% Engine failure -- fires at most once, ~0.2% probability per cycle
engine_failure_{id} :-
    race_started,
    \+ race_over,
    \+ engine_failure_{id}_fired,
    random(0, 1000, R), R < 2.

engine_failure_{id}I:>
    assert(engine_failure_{id}_fired),
    send_m(pitwall, send_message({id}_engine_failure, {id})).


% Push lap -- fires at most once per lap, ~10% probability
push_lap_{id} :-
    race_started,
    \+ race_over,
    \+ push_lap_{id}_fired,
    random(0, 100, R), R < 10.

push_lap_{id}I:>
    assert(push_lap_{id}_fired),
    send_m(pitwall, send_message({id}_push_lap, {id})).
\end{lstlisting}

\noindent
The \texttt{push\_lap\_\{id\}\_fired} flag is retracted at the start of each new lap
(\texttt{lap\_go} handler), so the event can fire once per lap.

\subsection{Non-deterministic Track Events}

The pit wall rolls \texttt{random\_track\_event/0} after each car completes a lap
(never on the last lap, at most once per lap via the \texttt{track\_event\_this\_lap} flag):

\begin{lstlisting}[language=Prolog, caption={random\_track\_event/0 in pitWallType}]
random_track_event :-
    if(track_event_this_lap,
       true,
       ( random(0, 10, R),
         if(R < 2,
            /* 20%: SAFETY CAR -- +10s to all cars */,
         if(R < 4,
            /* 20%: HEAVY RAIN -- +5s to all cars */,
            true)) )).
\end{lstlisting}

\subsection{Safety Car Chain}

The safety car involves a three-step message chain:

\begin{enumerate}[leftmargin=2em]
  \item Pit Wall detects \texttt{R < 2} $\Rightarrow$ asserts \texttt{track\_event\_this\_lap}, sends \texttt{deploy} to Safety Car and \texttt{safety\_car\_deployed} to all cars.
  \item At end-of-lap the pit wall sends \texttt{recall} to Safety Car $\Rightarrow$ Safety Car retracts \texttt{sc\_active} and sends \texttt{sc\_recalled} back to pit wall.
  \item Pit wall receives \texttt{sc\_recalled} $\Rightarrow$ prints ``GREEN FLAG!'' and sends \texttt{green\_flag} to all cars.
\end{enumerate}

% ================================================================
\section{How to Run}

\subsection{Native (WSL / Linux)}

\begin{lstlisting}[language=bash, numbers=none]
cd DALI/Examples/f1_race
bash startmas.sh
\end{lstlisting}

\texttt{startmas.sh} performs the following steps:
\begin{enumerate}[leftmargin=2em]
  \item Kills any stale SICStus / tmux processes on port 3010.
  \item Runs \texttt{generate\_agents.py} to sync DALI files with \texttt{agents.json}.
  \item Builds agent binaries and launches all agents inside a single \texttt{tmux}
        session named \texttt{f1\_race}.
\end{enumerate}

\subsection{Web Dashboard}

In a second terminal:

\begin{lstlisting}[language=bash, numbers=none]
cd DALI/Examples/f1_race
bash ui/run.sh
\end{lstlisting}

\noindent
\texttt{run.sh} creates a local Python venv (\texttt{ui/.venv}) on the first run and
installs Flask automatically. Subsequent launches skip \texttt{pip install} via a stamp
file (\texttt{.venv/.installed\_stamp}) and start in approximately 1\,second.

Open \url{http://localhost:5000} in a browser. The dashboard shows all agent panes
side-by-side with real-time auto-scrolling. If \texttt{agents.json} is modified while the
dashboard is open, it auto-detects the change within 5\,seconds and rebuilds the grid
without a page reload.

\subsection{Dashboard Controls}

\begin{longtable}{@{} l p{9cm} @{}}
\toprule
\textbf{UI element} & \textbf{Function} \\
\midrule
\endhead
\textbf{Restart MAS}    & Kills SICStus + tmux session, reruns \texttt{startmas.sh} \\
\textbf{Deploy SC}      & Sends deploy message to safety car immediately \\
\textbf{Recall SC}      & Recalls the safety car \\
Agent / Command bar     & Sends an arbitrary Prolog command to any agent pane \\
Pin button ($\downarrow$) & Toggles auto-scroll for that pane \\
Clear button ($\times$)  & Clears pane output \\
Minimise button ($-$)   & Minimises pane to tray \\
\bottomrule
\end{longtable}

% ================================================================
\section{Docker Deployment}

\subsection{Architecture}

The Docker Compose setup runs two containers that share a tmux socket volume, allowing
the dashboard to read agent panes with zero modifications to the agent logic.

\begin{lstlisting}[language={}, numbers=none, frame=none, backgroundcolor=\color{white}]
+----------------------------------+   +------------------------------+
|  mas  (ubuntu:22.04)             |   |  ui  (python:3.11-slim)      |
|  startmas.sh -> tmux f1_race     |   |  dashboard.py (Flask :5000)  |
|  SICStus mounted from host (ro)  |   |  reads tmux panes via volume |
+----------------------------------+   +------------------------------+
                 tmux_sock volume (/tmp/tmux-shared)
\end{lstlisting}

\noindent
\textbf{Startup order} (managed by Docker healthchecks):
\begin{enumerate}[leftmargin=2em]
  \item \texttt{ui} starts immediately and is healthy once \texttt{/api/config} responds.
  \item \texttt{mas} waits for \texttt{ui} to be healthy, then runs \texttt{startmas.sh}.
\end{enumerate}

\subsection{Quickstart}

\begin{lstlisting}[language=bash, numbers=none]
cd DALI/Examples/f1_race

# 1. Auto-detect SICStus and write .env  (run once)
bash docker/setup.sh

# 2. Build and start
docker compose up --build -d

# 3. Open dashboard
open http://localhost:5000
\end{lstlisting}

\subsection{Environment Variables}

\begin{longtable}{@{} l l p{7cm} @{}}
\toprule
\textbf{Variable} & \textbf{Default} & \textbf{Description} \\
\midrule
\endhead
\texttt{SICSTUS\_PATH} & \texttt{/usr/local/sicstus4.6.0} & Host path to SICStus install, mounted read-only into the \texttt{mas} container \\
\bottomrule
\end{longtable}

% ================================================================
\section{DALI Syntax Reference}

The following table summarises all DALI syntax constructs used in this project.

\begin{longtable}{@{} l p{5.5cm} p{5.5cm} @{}}
\toprule
\textbf{Syntax} & \textbf{Meaning} & \textbf{Example} \\
\midrule
\endhead
\texttt{nameE:>~Body.}         & React to external event \texttt{name} & \texttt{start\_raceE:> write('Go!').} \\[4pt]
\texttt{nameI:>~Body.}         & Fire internal event when \texttt{name} holds & \texttt{push\_lap\_ferrariI:> send\_m(...).} \\[4pt]
\texttt{name :- Cond.}         & Condition for internal event \texttt{nameI} & \texttt{push\_lap\_ferrari :- race\_started, ...} \\[4pt]
\texttt{messageA(ag,\,msg)}    & Send message (top-level clause) & \texttt{messageA(pitwall, lap\_done).} \\[4pt]
\texttt{send\_m(ag,\,msg)}     & Send message (safe inside \texttt{if/3}) & \texttt{send\_m(safety\_car, deploy).} \\[4pt]
\texttt{random(Lo,\,Hi,\,R)}   & Random integer $\texttt{Lo} \le R < \texttt{Hi}$ & \texttt{random(0, 10, R).} \\[4pt]
\texttt{if(Cond,\,Then,\,Else)} & Conditional & \texttt{if(R < 5, act, true).} \\[4pt]
\texttt{\textbackslash+~Goal}  & Negation-as-failure & \texttt{\textbackslash+ race\_over} \\[4pt]
\texttt{:-~Goal.}              & Directive (runs at load time) & \texttt{:- write('Agent ready!').} \\[4pt]
\texttt{keysort(+P,\,-S)}      & Sort \texttt{Key-Value} pairs by key & \texttt{keysort([3-b,\,1-a],\,S).} \\
\bottomrule
\end{longtable}

% ================================================================
\section{Shutdown}

\begin{lstlisting}[language=bash, numbers=none]
# Stop the MAS
tmux kill-session -t f1_race

# Stop the Flask dashboard
# Ctrl+C in the dashboard terminal

# Stop Docker
docker compose down
\end{lstlisting}

\end{document}
